\section{Generic BA reconstruction with repeated phases}
\label{app:BA-active-converse}

This appendix supplies the regularity step in
\cref{thm:BA-spectral-converse}. Throughout, $N=1$ and the two
phase lists are fixed, with arbitrary repetitions. A change of basis
is allowed only inside an equal-phase block; it then preserves both
the Hamiltonian and the fixed phase coefficients.

\subsection{Constant canonical bases and inactive directions}

For constant invertible block-diagonal matrices $\mathsf S_L,\mathsf S_M$ respecting
the phase lists, the coordinate changes
\begin{equation}\label{eq:BA-block-bases}
\begin{aligned}
 \boldsymbol\xi_L&\mapsto \mathsf S_L^{-\trans}\boldsymbol\xi_L,&
 \bar{\boldsymbol\xi}_L&\mapsto \mathsf S_L\bar{\boldsymbol\xi}_L,\\
 \boldsymbol\xi_M&\mapsto \mathsf S_M\boldsymbol\xi_M,&
 \bar{\boldsymbol\xi}_M&\mapsto \mathsf S_M^{-\trans}\bar{\boldsymbol\xi}_M,\\
 \boldsymbol\xi_{LM}&\mapsto \mathsf S_L\boldsymbol\xi_{LM}\mathsf S_M^{-1},&
 \bar{\boldsymbol\xi}_{LM}&\mapsto \mathsf S_L^{-\trans}\bar{\boldsymbol\xi}_{LM}\mathsf S_M^{\trans}
\end{aligned}
\end{equation}
are canonical and preserve \cref{eq:BA-H0-flat}. The corresponding
block moments transform by conjugation. Since they are conserved,
their diagonalizing bases can be chosen once at the initial time.

For a block of size $r$, let $n$ be the number of outside directions:
$n=M+1$ in the first group and $n=L+1$ in the second. Absorbing the
nonzero slope denominators into the incoming couplings writes its
moment as $\mathsf F\mathsf G$, with $\mathsf F$ of size $r\times n$ and $\mathsf G$ of size
$n\times r$. If $r>n$, then generically $\mathsf G\mathsf F$ is invertible and
\begin{equation}\label{eq:BA-inactive-splitting}
 \mathbb C^r=\operatorname{im}\mathsf F\oplus\ker\mathsf G
             =\operatorname{im}(\mathsf F\mathsf G)\oplus\ker(\mathsf F\mathsf G).
\end{equation}
The two subspaces are constant. In an adapted basis the $r-n$
directions in the second summand have zero incoming and outgoing
couplings. They stay zero and contribute only the spectral factor
$(\lambda-e_B-s_B\nu)^{r-n}$, where $e_B,s_B$ are their intercept
and slope. Delete them and repeat. The process terminates with
positive active dimensions $L_a,M_a$ such that every remaining block
has size at most its number of outside directions. Inverse basis
changes restore all original coordinates.

The reduction is generic: arbitrary maximal-rank coupling pairs can
be written as an embedded active pair, with zero extra directions,
followed by a block basis change. Thus the active coupling data are
not forced onto an additional exceptional algebraic subset.

\subsection{Regularity of the active spectral curve}

\begin{proposition}\label{prop:BA-active-regularity}
For every fixed phase profile, generic active initial modes have an
irreducible spectral curve, smooth in the affine plane, whose
normalization has distinct resolved infinity branches and nonspecial
scalar BA divisors. Each active moment block has simple nonzero
spectrum. The normalization genus is
\begin{equation}\label{eq:BA-active-genus}
 G_a=L_aM_a
 -\sum_{B\in\mathcal B_L}\frac{|B|(|B|-1)}2
 -\sum_{B\in\mathcal B_M}\frac{|B|(|B|-1)}2,
\end{equation}
where $\mathcal B_L,\mathcal B_M$ are the partitions into equal-phase
blocks of the active groups.
\end{proposition}

\begin{proof}
We check the affine curve, its connectivity, and its infinity data
separately. For a block of size $r\leq n$, the product $\mathsf F\mathsf G$ can have
arbitrary distinct nonzero eigenvalues: choose $\mathsf F$ to select $r$
outside directions and $\mathsf G$ to give a prescribed diagonal product.
Thus maximal coupling ranks and simple nonzero moment spectra are
nonempty open conditions. There are finitely many blocks, and their
conditions hold simultaneously on a nonempty open set of the
irreducible active mode space.

For affine smoothness put
$\mathsf M=\lambda\Id-\mathsf X-\nu\mathsf D$.
A right kernel vector supported in one original group would, by
the diagonal equations, be supported in one equal-phase block.
The incoming coupling map on that block is injective, so the vector
is zero. The analogous statement holds for left kernels. At a point
of corank $q$, put bases for the right and left kernels in the
columns of a $d\times q$ matrix $\mathsf R$ and the rows of a
$q\times d$ matrix $\mathsf S$, respectively. The derivative
transverse to the matrix rank stratum is
\[
 \delta\mathsf X\longmapsto \mathsf S\,\delta\mathsf X\,\mathsf R.
\]
This map is onto: a nonzero $q\times q$ annihilator $\mathsf Z$
would produce $\mathsf R\mathsf Z\mathsf S$
orthogonal to all off-diagonal-block variations, hence block diagonal.
A nonzero column of that block-diagonal matrix would be a forbidden
group-supported right kernel vector. The off-diagonal parameters
therefore meet every rank stratum transversely. Corank at least two
has codimension at least four, while $(\lambda,\nu)$ provide only two
variables. Such points are absent for generic parameters. The
corank-one determinant hypersurface is smooth, and generic smoothness
of its projection to parameter space gives smooth affine fibers
\cite[Chapter~III]{Hartshorne1977}.

For irreducibility, insert one index at a time, always keeping a
block's size at most its current outside dimension. Such an order
exists: if both groups had pending indices and neither insertion
were possible, the saturated blocks would imply both
$L_{\rm cur}\geq M_{\rm cur}+1$ and
$M_{\rm cur}\geq L_{\rm cur}+1$. If only one group has pending
indices, the final block-size bound allows its next insertion.
The $(1,1)$ base has an irreducible specialization
\[
 (\lambda-\rho)(\lambda-\nu-e_L)(\lambda+\nu-e_M)-\delta_0=0,
 \qquad \delta_0\ne0.
\]
As a quadratic in $\nu$ its discriminant has an odd pole at
$\lambda=\rho$, so is not a square.

For an insertion, at zero new outgoing coupling the characteristic
polynomial is $\ell F_{\rm old}$, with
$\ell=\lambda-e-s\nu$. Suppose $r$ old indices have the same phase
as the new index and $n$ directions lie outside its group. Eliminating
the other indices in its group gives
\[
 \deg_\nu F_{\rm old}(e+s\nu,\nu)=n-r.
\]
For $r>0$ the leading coefficient is a nonzero constant times the
old moment-block determinant; for $r=0$ it is the nonzero outside
leading determinant. The allowed insertion has $r+1\leq n$, so
the old curve meets the new line at a finite point. At that point
the smooth old pencil has corank one with nonzero outside components
in both kernels. Choose new coupling vectors $\boldsymbol a,\boldsymbol b$,
supported on the outside indices, so that
$\boldsymbol a\operatorname{adj}(\mathsf M_{\rm old})\boldsymbol b$ does not vanish there.
The perturbed determinant is
\[
 F_\varepsilon=\ell F_{\rm old}
       -\varepsilon \boldsymbol a\operatorname{adj}(\mathsf M_{\rm old})\boldsymbol b.
\]
A factorization deforming the coprime monic factors
$\ell,F_{\rm old}$ would force the first perturbation into
$(\ell,F_{\rm old})$, hence to vanish at every intersection.
This obstruction also excludes a factorization after finite algebraic
parameter extension. Indeed, put $\varepsilon=\tau^k$ and suppose
monic factors specialize to $\ell$ and $F_{\rm old}$. Write them as
\[
 G=\ell+\sum_{m\geq1}\tau^m a_m,\qquad
 J=F_{\rm old}+\sum_{m\geq1}\tau^m b_m,
 \qquad \deg_\lambda a_m<1,\quad
        \deg_\lambda b_m<\deg_\lambda F_{\rm old}.
\]
For $m<k$, induction on $m$ gives
$a_mF_{\rm old}+b_m\ell=0$. Coprimeness forces $\ell$ to divide
$a_m$, hence $a_m=b_m=0$ by the degree bounds. At order $k$, the
same coefficient equation would put
$-\boldsymbol a\operatorname{adj}(\mathsf M_{\rm old})\boldsymbol b$
in the ideal $(\ell,F_{\rm old})$, contrary to its nonzero value at
the selected intersection. Thus fractional parameter powers cannot
remove the obstruction. Induction proves generic geometric
irreducibility for every active phase profile.

Finally diagonalize the conserved moment blocks. Their simple
eigenvalues $\vartheta_i$ give the infinity expansions
\[
 \lambda=s_B\nu+e_B+\vartheta_i/\nu+O(\nu^{-2}),
 \qquad1\leq i\leq |B|.
\]
Branches in different phase blocks of one group have contact order
one, and those in the same block have contact order two. They are
distinct smooth points on the normalization. Counting these contacts
relative to the distinct-phase genus $L_aM_a$ gives
\cref{eq:BA-active-genus}. The singleton point remains smooth.

In the moment eigenbasis the leading spectral eigenvectors at infinity
form an invertible basis. The eigenline module is therefore saturated
on the normalization, including these tangent plane branches. Its
pushforward under $\nu$ is the original trivial bundle of rank
$L_a+M_a+1$. The line-bundle degree, last-component zeros and exact
sequence argument in \cref{thm:BA-spectral-converse} apply with genus
$G_a$, giving degree-$G_a$ nonspecial scalar divisors. The bilinear
kernel and differential \cref{eq:BA-converse-duality} have the same
divisor and puncture properties, so all reconstruction formulas apply.
\end{proof}

The original elementary coordinates are thus constant linear
combinations of the scalar BA evaluations in the active basis; the
pair matrix and its partner transform on both sides as in
\cref{eq:BA-block-bases}. This handles generic initial modes for
every fixed phase list. Additional singular spectral levels and
nonsemisimple moment specializations require separate limiting data.
